
\section{列国的总结}
\subsection{民數記 24:12-25 - 巴蘭預言以色列摩押以東亞瑪力基尼亞述}
巴蘭對巴勒說：「我豈不是對你所差遣到我那裡的使者說， 13 巴勒就是將他滿屋的金銀給我，我也不得越過耶和華的命，憑自己的心意行好行歹？耶和華說什麼，我就要說什麼。 14 現在我要回本族去。你來，我告訴你這民日後要怎樣待你的民。」 15 他就題起詩歌說：「比珥的兒子巴蘭說，眼目閉住的人說， 16 得聽神的言語、明白至高者的意旨、看見全能者的異象、眼目睜開而仆倒的人說： 17 我看他卻不在現時，我望他卻不在近日。有星要出於雅各，有杖要興於\textcolor{red}{以色列}，必打破\textcolor{red}{摩押}的四角，毀壞擾亂之子。 18 他必得\textcolor{red}{以東}為基業，又得仇敵之地\textcolor{red}{西珥}為產業，以色列必行事勇敢。 19 有一位出於雅各的，必掌大權，他要除滅城中的餘民。」 20 巴蘭觀看\textcolor{red}{亞瑪力}，就題起詩歌說：「亞瑪力原為諸國之首，但他終必沉淪。」 21 巴蘭觀看\textcolor{red}{基尼人}，就題起詩歌說：「你的住處本是堅固，你的窩巢做在巖穴中。 22 然而基尼必至衰微，直到\textcolor{red}{亞述}把你擄去。」 23 巴蘭又題起詩歌說：「哀哉！神行這事，誰能得活？ 24 必有人乘船從基提界而來，苦害\textcolor{red}{亞述}，苦害希伯，他也必至沉淪。」 25 於是巴蘭起來，回他本地去。巴勒也回去了。

\subsection{以賽亞書29:1-7/33:1-12}
\subsubsection{預言耶路撒冷必遭重災 (賽29:1-7)}
唉，亞利伊勒，亞利伊勒，大衛安營的城！任憑你年上加年，節期照常周流， 2 我終必使亞利伊勒困難，她必悲傷哀號，我卻仍以她為亞利伊勒。 3 我必四圍安營攻擊你，屯兵圍困你，築壘攻擊你。 4 你必敗落，從地中說話，你的言語必微細出於塵埃。你的聲音必像那交鬼者的聲音出於地，你的言語低低微微出於塵埃。

5 「你仇敵的群眾卻要像細塵，強暴人的群眾也要像飛糠。這事必頃刻之間，忽然臨到。 6 萬軍之耶和華必用雷轟、地震、大聲、旋風、暴風並吞滅的火焰，向她討罪。 7 那時，攻擊亞利伊勒\textcolor{red}{列國}的群眾，就是一切攻擊亞利伊勒和她的保障並使她困難的，必如夢景，如夜間的異象。 8 又必像飢餓的人，夢中吃飯，醒了仍覺腹空；或像口渴的人，夢中喝水，醒了仍覺發昏，心裡想喝。攻擊錫安山\textcolor{red}{列國}的群眾，也必如此。」

\subsubsection{耶和華在錫安被尊崇列國四散 (賽33:1-12)}
禍哉！你這毀滅人的，自己倒不被毀滅；行事詭詐的，人倒不以詭詐待你。你毀滅罷休了，自己必被毀滅；你行完了詭詐，人必以詭詐待你。 2 耶和華啊，求你施恩於我們，我們等候你！求你每早晨做我們的膀臂，遭難的時候為我們的拯救。 3 喧嚷的響聲一發，眾民奔逃；你一興起，\textcolor{red}{列國}四散。 4 你們所擄的必被斂盡，好像螞蚱吃(斂)盡禾稼；人要蹦在其上，好像蝗蟲一樣。 5 耶和華被尊崇，因他居在高處，他以公平、公義充滿錫安。 6 你一生一世必得安穩，有豐盛的救恩並智慧和知識，你以敬畏耶和華為至寶。

7 看哪，他們的豪傑在外頭哀號，求和的使臣痛痛哭泣。 8 大路荒涼，行人止息；敵人背約，藐視城邑，不顧人民。 9 地上悲哀衰殘，黎巴嫩羞愧枯乾，沙崙像曠野，巴珊和迦密的樹林凋殘。 10 耶和華說：「現在我要起來，我要興起，我要勃然而興！ 11 你們要懷的是糠秕，要生的是碎秸，你們的氣就是吞滅自己的火。 12 列邦必像已燒的石灰，像已割的荊棘，在火中焚燒。

\subsection{耶利米書}
\subsubsection{以怒杯喻列國刀劍之災 (25:15-31)}
\begin{table}[H]
\centering
\begin{tabular}{p{8.25cm}|p{9.25cm}}\hline\hline
\multicolumn{1}{c|}{耶利米}&\multicolumn{1}{c}{耶和華}\\\hline
&耶和華以色列的神對我如此說：「你從我手中接這杯憤怒的酒，使我所差遣你去的各國的民喝。 16 他們喝了就要東倒西歪，並要發狂，因我使刀劍臨到他們中間。」\\\hline
\multirow{2}{8.4cm}{17 我就從耶和華的手中接了這杯，給耶和華所差遣我去的各國的民喝， 18 就是\textcolor{red}{耶路撒冷}和\textcolor{red}{猶大}的城邑並耶路撒冷的君王與首領，使這城邑荒涼，令人驚駭、嗤笑、咒詛，正如今日一樣。 19 又有\textcolor{red}{埃及}王法老和他的臣僕、首領以及他的眾民， 20 並\textcolor{red}{雜族}的人民和\textcolor{red}{烏斯}地的諸王，與\textcolor{red}{非利士}地的諸王，\textcolor{red}{亞實基倫}，\textcolor{red}{加沙}，\textcolor{red}{以革倫}，\textcolor{red}{以及亞實突}剩下的人， 21 \textcolor{red}{以東}，\textcolor{red}{摩押}，\textcolor{red}{亞捫}人， 22 \textcolor{red}{推羅}的諸王，\textcolor{red}{西頓}的諸王，\textcolor{red}{海島}的諸王， 23\textcolor{red}{ 底但}，\textcolor{red}{提瑪}，\textcolor{red}{布斯}和一切剃周圍頭髮的， 24 \textcolor{red}{阿拉伯}的諸王，住\textcolor{red}{曠野雜族}人民的諸王， 25 \textcolor{red}{心利}的諸王，\textcolor{red}{以攔}的諸王，\textcolor{red}{瑪代}的諸王， 26 \textcolor{red}{北方遠近}的諸王以及天下\textcolor{red}{地上的萬國}喝了，以後\textcolor{red}{示沙克(就是巴比倫)}王也要喝。}
&27 「你要對他們說：『萬軍之耶和華以色列的神如此說：你們要喝，且要喝醉，要嘔吐，且要跌倒不得再起來，都因我使刀劍臨到你們中間。』 28 他們若不肯從你手接這杯喝，你就要對他們說：『萬軍之耶和華如此說：你們一定要喝！ 29 我既從稱為我名下的城起首施行災禍，你們能盡免刑罰嗎？你們必不能免，因為我要命刀劍臨到地上一切的居民。這是萬軍之耶和華說的。』\\\cline{2-2}
&「所以你要向他們預言這一切的話，攻擊他們，說：『耶和華必從高天吼叫，從聖所發聲，向自己的羊群大聲吼叫。他要向地上一切的居民呐喊，像踹葡萄的一樣。 31 必有響聲達到地極，因為耶和華與列國相爭。凡有血氣的，他必審問；至於惡人，他必交給刀劍。這是耶和華說的。』」\\\hline
\end{tabular}
\end{table}
 
\subsubsection{神憤怒的杯中純一不雜之災 (啟 14:8-12)}
又有第二位天使接著說：「\textcolor{red}{叫萬民喝邪淫、大怒之酒}的巴比倫大城傾倒了！傾倒了！」又有第三位天使接著他們大聲說：「若有人拜獸和獸像，在額上或在手上受了印記， 10 這人也\textcolor{red}{必喝神大怒的酒}，此酒斟在\textcolor{red}{神憤怒的杯}中純一不雜。他要在聖天使和羔羊面前，在火與硫磺之中受痛苦。 11 他受痛苦的煙往上冒，直到永永遠遠。那些拜獸和獸像，受牠名之印記的，晝夜不得安寧。」 12 聖徒的忍耐就在此，他們是守神誡命和耶穌真道的。


\subsubsection{以索與軛喻列國必侍巴比倫王,不服者必遭諸災 27:1-11}
\begin{table}[H]
\centering
\begin{tabular}{p{9.25cm}|p{8.25cm}}\hline\hline
\multicolumn{1}{c|}{列國都必服侍巴比倫王}&\multicolumn{1}{c}{不服者必遭諸災,服者在本地存留}\\\hline
\multirow{2}{9.4cm}{猶大王約西亞的兒子約雅敬(西底家,27:3)登基的時候，有這話從耶和華臨到耶利米說： 2 耶和華對我如此說：「你做繩索與軛，加在自己的頸項上， 3 藉那些來到耶路撒冷見猶大王西底家的使臣之手，把繩索與軛送到以東王、摩押王、亞捫王、推羅王、西頓王那裡。 4 且囑咐使臣傳於他們的主人說：『萬軍之耶和華以色列的神如此說： 5 我用大能和伸出來的膀臂創造大地和地上的人民、牲畜，我看給誰相宜，就把地給誰。 6 現在我將這些地都交給我僕人巴比倫王尼布甲尼撒的手，我也將田野的走獸給他使用。 7 列國都必服侍他和他的兒孫，直到他本國遭報的日期來到，那時多國和大君王要使他做他們的奴僕。}
&『無論哪一邦哪一國，不肯服侍這巴比倫王尼布甲尼撒，也不把頸項放在巴比倫王的軛下，我必用刀劍、饑荒、瘟疫刑罰那邦，直到我藉巴比倫王的手將他們毀滅。這是耶和華說的。 9 至於你們，不可聽從你們的先知和占卜的、圓夢的、觀兆的以及行邪術的，他們告訴你們說你們不致服侍巴比倫王。 10 他們向你們說假預言，要叫你們遷移，遠離本地，以致我將你們趕出去，使你們滅亡。\\\cline{2-2}
&11 但哪一邦肯把頸項放在巴比倫王的軛下服侍他，我必使那邦仍在本地存留，得以耕種居住。這是耶和華說的。』」\\\hline
\end{tabular}
\end{table}

\subsection{約珥書-在約沙法谷，審判四圍的列國 3:1-19}
到那日，我使猶大和耶路撒冷被擄之人歸回的時候， 2 我要聚集萬民，帶他們下到約沙法谷，在那裡施行審判。因為他們將我的百姓，就是我的產業以色列，分散在列國中，又分取我的地土， 3 且為我的百姓拈鬮，將童子換妓女，賣童女買酒喝。 4 推羅、西頓和非利士四境的人哪，你們與我何干？你們要報復我嗎？若報復我，我必使報應速速歸到你們的頭上。 5 你們既然奪取我的金銀，又將我可愛的寶物帶入你們宮殿（廟中）， 6 並將猶大人和耶路撒冷人賣給希臘人（雅完人），使他們遠離自己的境界， 7 我必激動他們離開你們所賣到之地，又必使報應歸到你們的頭上。 8 我必將你們的兒女賣在猶大人的手中，他們必賣給遠方示巴國的人。這是耶和華說的。

當在萬民中宣告說：要預備打仗，激動勇士，使一切戰士上前來！ 10 要將犁頭打成刀劍，將鐮刀打成戈矛。軟弱的要說：『我有勇力！』」 11 四圍的列國啊，你們要速速地來，一同聚集！耶和華啊，求你使你的大能者降臨！ 12 「萬民都當興起，上到約沙法谷，因為我必坐在那裡，審判四圍的列國。 13 開鐮吧！因為莊稼熟了。踐踏吧！因為酒榨滿了，酒池盈溢，他們的罪惡甚大。

14 「許多許多的人在斷定谷，因為耶和華的日子臨近斷定谷。 15 日月昏暗，星宿無光。 16 耶和華必從錫安吼叫，從耶路撒冷發聲，天地就震動。耶和華卻要做他百姓的避難所，做以色列人的保障。 17 你們就知道我是耶和華你們的神，且又住在錫安我的聖山。那時，耶路撒冷必成為聖，外邦人不再從其中經過。

18 「到那日，大山要滴甜酒，小山要流奶子，猶大溪河都有水流，必有泉源從耶和華的殿中流出來，滋潤什亭谷。 19 埃及必然荒涼，以東變為淒涼的曠野，都因向猶大人所行的強暴，又因在本地流無辜人的血。

但猶大必存到永遠，耶路撒冷必存到萬代。 21 我未曾報復（洗除）流血的罪，現在我要報復，因為耶和華住在錫安。

\subsection{撒迦利亞書-耶和華用災殃攻擊與耶路撒冷爭戰的列國人 14:12-21}
耶和華用災殃攻擊那與耶路撒冷爭戰的列國人，必是這樣：他們兩腳站立的時候，肉必消沒，眼在眶中乾癟，舌在口中潰爛。 13 那日，耶和華必使他們大大擾亂，他們各人彼此揪住，舉手攻擊。 14 猶大也必在耶路撒冷爭戰。那時四圍各國的財物，就是許多金銀、衣服，必被收聚。 15 那臨到馬匹、騾子、駱駝、驢和營中一切牲畜的災殃是與那災殃一般。

所有來攻擊耶路撒冷列國中剩下的人，必年年上來敬拜大君王萬軍之耶和華，並守住棚節。 17 地上萬族中，凡不上耶路撒冷敬拜大君王萬軍之耶和華的，必無雨降在他們的地上。 18 埃及族若不上來，雨也不降在他們的地上。凡不上來守住棚節的列國人，耶和華也必用這災攻擊他們。 19 這就是埃及的刑罰和那不上來守住棚節之列國的刑罰。 20 當那日，馬的鈴鐺上必有「歸耶和華為聖的」這句話。耶和華殿內的鍋必如祭壇前的碗一樣。 21 凡耶路撒冷和猶大的鍋都必歸萬軍之耶和華為聖，凡獻祭的都必來取這鍋，煮肉在其中。當那日，在萬軍之耶和華的殿中必不再有迦南人。

\subsection{撒迦利亞書-耶和華用災殃攻擊與耶路撒冷爭戰的列國人 14:12-21}





\section{亞述}
\subsection{亞述掳走以色列}
\subsubsection{以色列被掳亞述 (何9:1-9)}

以色列啊，不要像外邦人歡喜快樂，因為你行邪淫離棄你的神，在各穀場上如妓女喜愛賞賜。 2 穀場和酒榨都不夠以色列人使用，新酒也必缺乏。 3 他們必不得住耶和華的地，以法蓮卻要\textcolor{red}{歸回埃及}，必在\textcolor{red}{亞述吃不潔淨的食物}。 4 他們必不得向耶和華奠酒，即便奠酒也不蒙悅納。他們的祭物必如居喪者的食物，凡吃的必被玷汙，因他們的食物只為自己的口腹，必不奉入耶和華的殿。 5 在大會的日子，到耶和華的節期，你們怎樣行呢？ 6 看哪，他們逃避災難，\textcolor{red}{埃及人}必收殮他們的屍首，\textcolor{red}{摩弗人}必葬埋他們的骸骨。他們用銀子做的美物上必長蒺藜，他們的帳篷中必生荊棘。 7 以色列人必知道降罰的日子臨近，報應的時候來到。民說：「做先知的是愚昧，受靈感的是狂妄。」皆因他們多多作孽，大懷怨恨。 8 以法蓮曾做我神守望的，至於先知，在他一切的道上作為捕鳥人的網羅，在他神的家中懷怨恨。 9 以法蓮深深地敗壞，如在基比亞的日子一樣。耶和華必記念他們的罪孽，追討他們的罪惡。

%\subsection{阿摩司書}
\subsubsection{擄到大馬士革以外 (摩5:25-27)}
以色列家啊，你們在曠野四十年，豈是將祭物和供物獻給我呢？ 26 你們抬著為自己所造之摩洛的帳幕和偶像的龕，並你們的神星。 27 所以，我要把你們\textcolor{red}{擄到大馬士革以外}。」這是耶和華，名為萬軍之神說的。

\subsubsection{耶和華必興一國為以色列之敵 (摩6:1-14)}
國為列國之首，人最著名，且為以色列家所歸向，在錫安和撒馬利亞山安逸無慮的，有禍了！ 2 你們要過到甲尼察看，從那裡往大城哈馬去，又下到非利士人的迦特，看那些國比你們的國還強嗎？境界比你們的境界還寬嗎？ 3 你們以為降禍的日子還遠，坐在位上盡行強暴 a。 4 你們躺臥在象牙床上，舒身在榻上，吃群中的羊羔、棚裡的牛犢。 5 彈琴鼓瑟唱消閒的歌曲，為自己製造樂器，如同大衛所造的。 6 以大碗喝酒，用上等的油抹身，卻不為約瑟的苦難擔憂。

7 所以這些人必在被擄的人中首先被擄，舒身的人荒宴之樂必消滅了。 8 主耶和華萬軍之神指著自己起誓說：「我憎惡雅各的榮華，厭棄他的宮殿，因此我必將城和其中所有的都交付敵人。」 9 那時，若在一房之內剩下十個人，也都必死。 10 死人的伯叔，就是燒他屍首的，要將這屍首搬到房外，問房屋內間的人說：「你那裡還有人沒有？」他必說：「沒有。」又說：「不要作聲！因為我們不可提耶和華的名。」 11 看哪，耶和華出令，大房就被攻破，小屋就被打裂。

12 馬豈能在崖石上奔跑？人豈能在那裡用牛耕種呢？你們卻使公平變為苦膽，使公義的果子變為茵陳。 13 你們喜愛虛浮的事，自誇說：「我們不是憑自己的力量取了角嗎？」

14 耶和華萬軍之神說：「以色列家啊，我必興起一國攻擊你們，他們必欺壓你們，從哈馬口直到亞拉巴的河。」

\subsubsection{亞述為主之杖，心裡倒想毀滅 (賽10:5-11)}
\begin{table}[H]
\arrayrulecolor{black}
\centering 
\begin{tabular}{p{5.5cm}|p{11.25cm}}\hline\hline
亞述為主之杖&心裡倒想毀滅\\\hline
亞述是我怒氣的棍，手中拿我惱恨的杖。我要打發他攻擊褻瀆的國民，吩咐他攻擊我所惱怒的百姓，搶財為擄物，奪貨為掠物，將他們踐踏，像街上的泥土一樣。
&然而，他不是這樣的意思，他心也不這樣打算，他心裡倒想毀滅，剪除不少的國。他說：『我的臣僕豈不都是王嗎？迦勒挪豈不像迦基米施嗎？哈馬豈不像亞珥拔嗎？撒馬利亞豈不像大馬士革嗎？我手已經夠到有偶像的國，這些國雕刻的偶像過於耶路撒冷和撒馬利亞的偶像。我怎樣待撒馬利亞和其中的偶像，豈不照樣待耶路撒冷和其中的偶像嗎？』\\\hline
\end{tabular}
\end{table}

\subsubsection{亞述攻打錫安之路线 (賽10:24-34)}
\begin{table}[H]
\arrayrulecolor{black}
\centering 
\begin{tabular}{p{8.5cm}|p{5.75cm}|p{2.75cm}}\hline\hline
亞述的重擔必離開錫安&亞述王使万民惊恐&神以驚嚇削樹枝\\\hline
所以主萬軍之耶和華如此說：「住錫安我的百姓啊，亞述王雖然用棍擊打你，又照埃及的樣子舉杖攻擊你，你卻不要怕他。因為還有一點點時候，向你們發的憤恨就要完畢，我的怒氣要向他發作，使他滅亡。」萬軍之耶和華要興起鞭來攻擊他，好像在俄立磐石那裡殺戮米甸人一樣；耶和華的杖要向海伸出，把杖舉起，像在埃及一樣。到那日，亞述王的重擔必離開你的肩頭，他的軛必離開你的頸項，那軛也必因肥壯的緣故撐斷 (因膏油的緣故毀壞)。
&亞述王來到亞葉，經過米磯崙，在密抹安放輜重。他們過了隘口，在迦巴住宿。拉瑪人戰兢，掃羅的基比亞人逃跑。迦琳的居民(女子)哪，要高聲呼喊！萊煞人哪，須聽！哀哉，困苦的亞拿突啊！瑪得米那人躲避，基柄的居民逃遁。當那日，亞述王要在挪伯歇兵，向錫安女子的山，就是耶路撒冷的山，掄手攻他。&看哪，主萬軍之耶和華以驚嚇削去樹枝，長高的必被砍下，高大的必被伐倒。稠密的樹林，他要用鐵器砍下，黎巴嫩的樹木必被大能者伐倒。\\\hline
\end{tabular}
\end{table}



\subsubsection{使驚動靈進入亞述王心,在本地倒在刀下,亞述王因驕致禍 (賽37:5-7/37:21-32)}
希西家王的臣僕就去見以賽亞。 6 以賽亞對他們說：「要這樣對你們的主人說：『耶和華如此說：你聽見亞述王的僕人褻瀆我的話，不要懼怕。 7 我必驚動(使靈進入)他的心，他要聽見風聲，就歸回本地，我必使他在那裡倒在刀下。』」

%\subsubsection{以賽亞預言亞述王因驕致禍 37:21-32}
亞摩斯的兒子以賽亞就打發人去見希西家，說：「耶和華以色列的神如此說：你既然求我攻擊亞述王西拿基立， 22 所以耶和華論他這樣說：『錫安的處女藐視你，嗤笑你，耶路撒冷的女子向你搖頭。 23 你辱罵誰，褻瀆誰？揚起聲來，高舉眼目攻擊誰呢？乃是攻擊以色列的聖者！ 24 你藉你的臣僕辱罵主，說：「我率領許多戰車上山頂，到黎巴嫩極深之處，我要砍伐其中高大的香柏樹和佳美的松樹，我必上極高之處，進入肥田的樹林。 25 我已經挖井喝水，我必用腳掌踏乾埃及的一切河。」』

26 「耶和華說：『你豈沒有聽見我早先所做的，古時所立的嗎？現在藉你使堅固城荒廢，變為亂堆。 27 所以其中的居民力量甚小，驚惶羞愧，他們像野草，像青菜，如房頂上的草，又如田間未長成的禾稼。 28 你坐下，你出去，你進來，你向我發烈怒，我都知道。 29 因你向我發烈怒，又因你狂傲的話達到我耳中，我就要用鉤子鉤上你的鼻子，把嚼環放在你口裡，使你從原路轉回去。』

以色列人哪，我賜你們一個證據：你們今年要吃自生的，明年也要吃自長的，至於後年，你們要耕種收割，栽植葡萄園，吃其中的果子。 31 猶大家所逃脫餘剩的，仍要往下扎根，向上結果。 32 必有餘剩的民從耶路撒冷而出，必有逃脫的人從錫安山而來。萬軍之耶和華的熱心必成就這事。
「所以耶和華論亞述王如此說：他必不得來到這城，也不在這裡射箭，不得拿盾牌到城前，也不築壘攻城。 34 他從哪條路來，必從哪條路回去，必不得來到這城。這是耶和華說的。 35 因我為自己的緣故，又為我僕人大衛的緣故，必保護拯救這城。」






\subsection{亞述结局}
\subsubsection{耶和華必罰亞述的自大，將其焚燒淨盡 (賽10:12-19)}
\begin{table}[H]
\arrayrulecolor{black}
\centering 
\begin{tabular}{p{8.5cm}|p{8.75cm}}\hline\hline
耶和華必罰亞述的自大&將其焚燒淨盡\\\hline
主在錫安山和耶路撒冷成就他一切工作的時候，主說：「我必罰亞述王自大的心和他高傲眼目的榮耀。
&\multirow{2}{8.9cm}{斧豈可向用斧砍木的自誇呢？鋸豈可向用鋸的自大呢？好比棍掄起那舉棍的，好比杖舉起那非木的人。因此主萬軍之耶和華必使亞述王的肥壯人變為瘦弱，在他的榮華之下必有火著起如同焚燒一樣。以色列的光必如火，他的聖者必如火焰，在一日之間，將亞述王的荊棘和蒺藜焚燒淨盡。又將他樹林和肥田的榮耀全然燒盡，好像拿軍旗的昏過去一樣。他林中剩下的樹必稀少，就是孩子也能寫其數。}\\\cline{1-1}
因為他說：『我所成就的事是靠我手的能力和我的智慧，我本有聰明。我挪移列國的地界，搶奪他們所積蓄的財寶，並且我像勇士使坐寶座的降為卑。我的手夠到列國的財寶，好像人夠到鳥窩；我也得了全地好像人拾起所棄的雀蛋.沒有動翅膀的.沒有張嘴的.也沒有鳴叫的。』&\\\hline
\end{tabular}
\end{table}

\subsubsection{打折亞述人，陀斐特杖擊火烧，避刀的少年人成為服苦 (賽14:24-27/30:31-33/31:8-9)}
萬軍之耶和華起誓說：「我怎樣思想，必照樣成就；我怎樣定意，必照樣成立。\textcolor{blue}{ 就是在我地上打折亞述人，在我山上將他踐踏。他加的軛必離開以色列人，他加的重擔必離開他們的肩頭。」} 26 這是向全地所定的旨意，這是向萬國所伸出的手。 27 萬軍之耶和華既然定意，誰能廢棄呢？他的手已經伸出，誰能轉回呢？
%\subsubsection{杖擊亞述人,在陀斐特火與許多木柴為王預備 30:31-33}
\textcolor{blue}{亞述人必因耶和華的聲音驚惶，耶和華必用杖擊打他。 32 耶和華必將命定的杖加在他身上，每打一下，人必擊鼓彈琴。打仗的時候，耶和華必掄起手來與他交戰。} 33 原來陀斐特又深又寬，早已為王預備好了，其中堆的是火與許多木柴，耶和華的氣如一股硫磺火，使它著起來。
%\subsubsection{亞述人倒在非人的刀下,少年人逃避這刀的必成為服苦的,首領因大旗驚惶 31:8-9}
「亞述人必倒在刀下，並非人的刀；有刀要將他吞滅，並非人的刀。他必逃避這刀，他的少年人必成為服苦的。 9 他的磐石必因驚嚇挪去，他的首領必因大旗驚惶。」這是那有火在錫安、有爐在耶路撒冷的耶和華說的。

\subsubsection{亞述擄基尼，有船從基提界來苦害亞述,希伯，他也必至沉淪 (民24:21-25)}
巴蘭觀看基尼人，就題起詩歌說：「你的住處本是堅固，你的窩巢做在巖穴中。 22 然而基尼必至衰微，直到亞述把你擄去。」 23 巴蘭又題起詩歌說：「哀哉！神行這事，誰能得活？ 24 必有人乘船從基提界而來，苦害亞述，苦害希伯，他也必至沉淪。」 25 於是巴蘭起來，回他本地去。巴勒也回去了。











\subsection{尼尼微的默示}
\subsubsection{耶和華以烈怒施諸敵以恩慈待己兵 (鴻1:1-15)}
\begin{table}[H]
\centering
\begin{tabular}{p{10.25cm}|p{7.cm}}\hline\hline
\multicolumn{1}{c|}{尼尼微}&\multicolumn{1}{c}{猶大}\\\hline
\multirow{4}{10.5cm}{論尼尼微的默示,就是伊勒歌斯人那鴻所得的默示.耶和華是忌邪施報的神,耶和華施報大有憤怒,向他的敵人施報,向他的仇敵懷怒。耶和華不輕易發怒,大有能力,萬不以有罪的為無罪.他乘\textcolor{red}{旋風和暴風}而來,雲彩為他腳下的塵土.他斥責海,使海乾了,使一切\textcolor{red}{江河乾涸}.巴珊和迦密的樹林衰殘,黎巴嫩的花草也衰殘了.大\textcolor{red}{山因他震動},小山也都消化.大\textcolor{red}{地在他面前突起},世界和住在其間的也都如此.他發憤恨,誰能立得住呢?他發烈怒,誰能當得起呢?他的憤怒如火傾倒,磐石因他崩裂.耶和華本為善,在患難的日子為人的保障,並且認得那些投靠他的人.但他必\textcolor{red}{以漲溢的洪水淹沒尼尼微，又驅逐仇敵進入黑暗}.尼尼微人哪,設何謀攻擊耶和華呢?他必將你們滅絕淨盡,災難不再興起.你們像叢雜的荊棘,像喝醉了的人,又如枯乾的碎秸\textcolor{red}{全然燒滅}.}
&有一人從你那裡出來，圖謀邪惡，設惡計攻擊耶和華。耶和華如此說：\textcolor{red}{尼尼微}雖然勢力充足，人數繁多，也被剪除，歸於無有。\\\cline{2-2}
&\textcolor{red}{猶大}啊,我雖然使你受苦,卻不再使你受苦.現在我必從你頸項上折斷他的軛,扭開他的繩索\\\cline{2-2}
&耶和華已經出令，指著\textcolor{red}{尼尼微}說：「你名下的人必不留後。我必從你神的廟中除滅雕刻的偶像和鑄造的偶像，我必因你鄙陋使你歸於墳墓。」\\\cline{2-2}
&看哪,有報好信傳平安之人的腳登山說:「\textcolor{red}{猶大}啊,可以守你的節期,還你所許的願吧!因為那惡人不再從你中間經過,他已滅絕淨盡了」\\\hline
\end{tabular}
\end{table}


\subsubsection{以其得勝之軍傾覆尼尼微 (鴻2:1-13)}
尼尼微啊，那打碎邦國的上來攻擊你，你要看守保障，謹防道路，使腰強壯，大大勉力。 

耶和華復興雅各的榮華，好像以色列的榮華一樣，因為使地空虛的已經使雅各和以色列空虛，將他們的葡萄枝毀壞了。 3 他勇士的盾牌是紅的，精兵都穿朱紅衣服。在他預備爭戰的日子，戰車上的鋼鐵閃爍如火，柏木把的槍也掄起來了。 4 車輛在街上(城外)急行，在寬闊處奔來奔去，形狀如火把，飛跑如閃電。 5 尼尼微王招聚他的貴胄，他們步行絆跌，速上城牆，預備擋牌。 6 河閘開放，宮殿沖沒。 7 王后蒙羞，被人擄去，宮女捶胸，哀鳴如鴿。此乃命定之事。

8 尼尼微自古以來充滿人民，如同聚水的池子，現在居民卻都逃跑。雖有人呼喊說「站住！站住！」，卻無人回顧。 9 你們搶掠金銀吧！因為所積蓄的無窮，華美的寶器無數。 10 尼尼微現在空虛荒涼，人心消化，雙膝相碰，腰都疼痛，臉都變色。 11 獅子的洞和少壯獅子餵養之處在哪裡呢？公獅、母獅、小獅遊行，無人驚嚇之地在哪裡呢？ 12 公獅為小獅撕碎許多食物，為母獅掐死活物，把撕碎的、掐死的充滿牠的洞穴。 13 萬軍之耶和華說：「我與你為敵，必將你的車輛焚燒成煙，刀劍也必吞滅你的少壯獅子。我必從地上除滅你所撕碎的，你使者的聲音必不再聽見。」

\subsubsection{尼尼微因惡受罰 (鴻3:1-19)}
禍哉，這流人血的城！充滿謊詐和強暴，搶奪的事總不止息。 2 鞭聲響亮，車輪轟轟，馬匹踢跳，車輛奔騰， 3 馬兵爭先，刀劍發光，槍矛閃爍，被殺的甚多，屍首成了大堆，屍骸無數，人碰著而跌倒； 4 都因那美貌的妓女多有淫行，慣行邪術，藉淫行誘惑列國，用邪術誘惑(賣)多族。 5 萬軍之耶和華說：「我與你為敵。我必揭起你的衣襟，蒙在你臉上，使列國看見你的赤體，使列邦觀看你的醜陋。 6 我必將可憎汙穢之物拋在你身上，辱沒你，為眾目所觀。 7 凡看見你的，都必逃跑離開你，說：『尼尼微荒涼了！』有誰為你悲傷呢？我何處尋得安慰你的人呢？」

8 你豈比挪亞們強呢？挪亞們坐落在眾河之間，周圍有水，海 (尼羅河)做她的濠溝，又做她的城牆。 9 古實和埃及是她無窮的力量，弗人和路比族是她的幫手。 10 但她被遷移，被擄去，她的嬰孩在各市口上也被摔死，人為她的尊貴人拈鬮，她所有的大人都被鏈子鎖著。 11 你也必喝醉，必被埋藏，並因仇敵的緣故尋求避難所。 12 你一切保障必像無花果樹上初熟的無花果，若一搖撼，就落在想吃之人的口中。 13 你地上的人民如同婦女，你國中的關口向仇敵敞開，你的門閂被火焚燒。 14 你要打水預備受困，要堅固你的保障，踹土和泥修補磚窯。 15 在那裡，火必燒滅你，刀必殺戮你，吞滅你如同蝻子。任你加增人數多如蝻子，多如蝗蟲吧！ 16 你增添商賈多過天上的星，蝻子吃盡而去。 17 你的首領多如蝗蟲，你的軍長彷彿成群的螞蚱，天涼的時候齊落在籬笆上，日頭一出便都飛去，人不知道落在何處。 18 亞述王啊，你的牧人睡覺，你的貴胄安歇，你的人民散在山間無人招聚。 19 你的損傷無法醫治，你的傷痕極其重大。凡聽你信息的，必都因此向你拍掌，你所行的惡誰沒有時常遭遇呢？


\subsubsection{耶和華必伸手攻擊亞述尼尼微荒涼乾旱 (番2:13-15)}
耶和華必伸手攻擊北方，毀滅\textcolor{red}{亞述}，使\textcolor{red}{尼尼微}荒涼，又乾旱如曠野。 14 群畜，就是各國(類)的走獸必臥在其中，鵜鶘和箭豬要宿在柱頂上，在窗戶內有鳴叫的聲音，門檻都必毀壞，香柏木已經露出。 15 這是素來歡樂安然居住的城，心裡說「唯有我，除我以外再沒有別的」，現在何竟荒涼，成為野獸躺臥之處！凡經過的人，都必搖手嗤笑她。




\subsection{以色列犹大结局}
\subsubsection{以色列和亞述关系  (何12:13-16/14:12-13)}
以法蓮吃風，且追趕東風，時常增添虛謊和強暴，與亞述立約，把油送到埃及。 2 耶和華與猶大爭辯，必照雅各所行的懲罰他，按他所做的報應他。


\subsubsection{為亞述餘剩的百姓,必有一條大道歸聖山過去不致濕腳  (賽11:13-16/27:12-13)}
\begin{table}[H]
\centering
\begin{tabular}{p{10.5cm}|p{6.0cm}}\hline\hline
\multicolumn{1}{c|}{大河必有一條大道歸聖山過去不致濕腳11:13-16}&\multicolumn{1}{c}{被趕散的都要來在聖山敬拜27:12-13}\\\hline
以法蓮的嫉妒就必消散，擾害猶大的必被剪除；以法蓮必不嫉妒猶大，猶大也不擾害以法蓮。 14 他們要向西飛，撲在非利士人的肩頭上 (西界)，一同擄掠東方人，伸手按住以東和摩押，亞捫人也必順服他們。 15 耶和華必使埃及海汊枯乾，\textcolor{blue}{掄手用暴熱的風使大河分為七條，令人過去不致濕腳。為主餘剩的百姓，就是從亞述剩下回來的，必有一條大道，如當日以色列從埃及地上來一樣}。
&以色列人哪，到那日，耶和華必從\textcolor{blue}{大河直到埃及小河，將你們一一地收集，如同人打樹拾果一樣。當那日，必大發角聲，在亞述地將要滅亡的，並在埃及地被趕散的都要來，他們就在耶路撒冷聖山上敬拜耶和華。}\\\hline
\end{tabular}
\end{table}

\subsection{其他关于亚述的经文}
\subsubsection{亚述结局的经文-賽何摩鴻番}
\begin{itemize}
\item 賽10:5-19,24-34/11:13-16/14:24-27 /27:12-13/30:31-33/31:8-9/37:5-7/37:21-32
\item 何 9:1-9
\item 摩5:25-27/6:1-14
\item 鴻 1:1-15/2:1-13/3:1-19
\item 番 2:13-15
\end{itemize}

\subsubsection{地理位置 （創世記 2:14）}

有河從伊甸流出來，滋潤那園子，從那裡分為四道。 11 第一道名叫比遜，就是環繞哈腓拉全地的。在那裡有金子， 12 並且那地的金子是好的，在那裡又有珍珠和紅瑪瑙。 13 第二道河名叫基訓，就是環繞古實全地的。 14 第三道河名叫底格里斯，流在亞述的東邊。第四道河就是幼發拉底河。 15 耶和華神將那人安置在伊甸園，使他修理看守。

\subsubsection{家谱，人种}
創世記10:12/22
含的兒子是古實、麥西、弗、迦南。 7 古實的兒子是西巴、哈腓拉、撒弗他、拉瑪、撒弗提迦。拉瑪的兒子是示巴、底但。 8 古實又生寧錄，他為世上英雄之首。 9 他在耶和華面前是個英勇的獵戶，所以俗語說：「像寧錄在耶和華面前是個英勇的獵戶。」 10 他國的起頭是巴別、以力、亞甲、甲尼，都在\textcolor{red}{示拿地}。 11 他從那地出來往亞述去，建造尼尼微、利河伯、迦拉， 12 和尼尼微、迦拉中間的利鮮，這就是那大城。 13 麥西生路低人、亞拿米人、利哈比人、拿弗土希人、 14 帕斯魯細人、迦斯路希人、迦斐托人，從迦斐托出來的有非利士人。

\textcolor{red}{創世記10:21-31}
雅弗的哥哥閃，是希伯子孫之祖，他也生了兒子。 22 閃的兒子是以攔、\textcolor{red}{亞述}、亞法撒、路德、亞蘭。 23 亞蘭的兒子是烏斯、戶勒、基帖、瑪施。 24 亞法撒生沙拉，沙拉生希伯。 25 希伯生了兩個兒子，一個名叫法勒 a，因為那時人就分地居住；法勒的兄弟名叫約坍。 26 約坍生亞摩答、沙列、哈薩瑪非、耶拉、 27 哈多蘭、烏薩、德拉、 28 俄巴路、亞比瑪利、示巴、 29 阿斐、哈腓拉、約巴，這都是約坍的兒子。 30 他們所住的地方，是從米沙直到西發東邊的山。 31 這就是閃的子孫，各隨他們的宗族、方言、所住的地土、邦國。

\textcolor{red}{歷代志上 1:17-27}
閃的兒子是以攔、亞述、亞法撒、路德、亞蘭、烏斯、戶勒、基帖、米設 c。 18 亞法撒生沙拉，沙拉生希伯。 19 希伯生了兩個兒子，一個名叫法勒 d，因為那時人就分地居住；法勒的兄弟名叫約坍。 20 約坍生亞摩答、沙列、哈薩瑪非、耶拉、 21 哈多蘭、烏薩、德拉、 22 以巴錄、亞比瑪利、示巴、 23 阿斐、哈腓拉、約巴，這都是約坍的兒子。24 閃生亞法撒，亞法撒生沙拉， 25 沙拉生希伯，希伯生法勒，法勒生拉吳， 26 拉吳生西鹿，西鹿生拿鶴，拿鶴生他拉， 27 他拉生亞伯蘭，亞伯蘭就是亞伯拉罕。

\textcolor{red}{創世記25:18}
撒拉的使女埃及人夏甲給亞伯拉罕所生的兒子是以實瑪利。 13 以實瑪利兒子們的名字，按著他們的家譜記在下面。以實瑪利的長子是尼拜約，又有基達、亞德別、米比衫、 14 米施瑪、度瑪、瑪撒、 15 哈大、提瑪、伊突、拿非施、基底瑪。 16 這是以實瑪利眾子的名字，照著他們的村莊、營寨，做了十二族的族長。 17 以實瑪利享壽一百三十七歲，氣絕而死，歸到他列祖 b那裡。 18 他子孫的住處在他眾弟兄東邊，從哈腓拉直到埃及前的書珥，正在亞述的道上。

https://cnbible.com/hebrew/804.htm
Hosea 12:2; Hosea 14:4; Isaiah 10:5; Isaiah 14:25; Isaiah 19:23 (twice in verse); Isaiah 19:24,25; Isaiah 23:13; Isaiah 30:31; Isaiah 31:8; Isaiah 52:4; Lamentations 5:6; Ezekiel 23:5; Ezekiel 27:23; Ezekiel 32:22 (here feminine) Zechariah 10:11; Psalm 83:9 perhaps read גְּשׁוּר, compare 2 Samuel 2:9 below אֲשׁוּרִי; or (if Psalm 83 be late) regard אַשּׁוּר (like עֲמָלֵק ib.) as used because of ancient significance; sometimes personified as one Isaiah 10:5; Ezekiel 31:3 (but strike out Co q. v.), compare also Micah 5:4; Micah 5:5; Zephaniah 2:13; ׳מַחֲנֵה א 2 Kings 19:32 = Isaiah 37:36; ׳בְּנֵי א Ezekiel 16:28; Ezekiel 23:7,9,12,23.

3 land of Assyria Genesis 2:14; Genesis 10:11; Hosea 5:13; Hosea 7:11; Hosea 8:9; Hosea 9:3; Hosea 10:6; Isaiah 11:11,16; Isaiah 19:23; Jeremiah 2:18,36; Micah 7:12; Zechariah 10:10; אַשּׁ֫וּרָה Genesis 25:18; Isaiah 19:23; 2 Kings 15:29; 2 Kings 17:6,23; 2 Kings 18:11; אֶרֶץ אַשּׁוּר Isaiah 7:18; Isaiah 27:13; Hosea 11:11; Micah 5:5.

4 especially מֶלֶךְ אַשּׁוּר Isaiah 8:4; Isaiah 10:12; Isaiah 20:1,4,6 (probably gloss Isaiah 7:17,20; Isaiah 8:7) 2 Kings 15:19 41t. 2Kings; 14t. Isaiah 36-38; 1 Chronicles 5:6 (אַשֻּׁר) + 13 t. Chronicles; also Jeremiah 50:17,18; Nahum 3:18; Ezra 4:2; (only Ezra 6:22 of Persian or any king not strictly Assyrian); note also ׳הַמֶּלֶךְ א Isaiah 36:8,16 (׳א perhaps gloss, compare Di who holds same view as to 2 Kings 18:23,31); ׳מַלְכֵי א 2 Kings 19:11,17 = Isaiah 37:11,18; 2Chronicles 28:16; 30:6; Nehemiah 9:32.









